\documentclass{article}

\usepackage{iclr2026_conference, times}

\usepackage{hyperref}
\usepackage{url}
\usepackage{booktabs}
\usepackage{graphicx}
\usepackage{amsmath}
\usepackage{microtype}
\usepackage{float}

\title{Tiny Reproduction: Validating Magnitude-Based Neural Network Pruning\\
\large A Small-Scale Study of One-Shot vs.\ Iterative Pruning on MNIST}

\author{%
  Sebastian Bohrt \\
  ECE 57000, Purdue University \\
}

\newcommand{\fix}{\marginpar{FIX}}
\newcommand{\new}{\marginpar{NEW}}
\iclrfinalcopy

\begin{document}

\maketitle

\begin{abstract}
Large neural networks are known to be heavily over-parameterized, yet this redundancy
is commonly exploited after training to reduce model size through weight pruning.
Han et al.\ \citeyear{han2015learning} showed that magnitude-based pruning can remove
the majority of weights with negligible accuracy loss, provided that fine-tuning is applied
after pruning.
We reproduce the core claim of that work at small scale on MNIST using a three-layer
multilayer perceptron (MLP).
Specifically, we compare \emph{one-shot} global magnitude pruning against
\emph{iterative} pruning with inter-round fine-tuning across sparsity levels from 0\%
to 95\%.
Our results confirm the original trend with striking clarity: accuracy remains
near-baseline under both methods up to ${\sim}$70\% sparsity, but one-shot pruning
collapses sharply beyond 80\%---reaching only 56.67\% at 95\% sparsity---while iterative
pruning retains 96.51\%, a 39.8 percentage-point gap that validates fine-tuning as the
critical recovery mechanism.
We additionally find that naive layer-wise one-shot pruning performs catastrophically
worse than global pruning (42.88\% vs.\ 80.43\% at 90\%), revealing a key failure mode
when small output layers are over-compressed.
Code is available at \url{https://github.com/sbohrt/ece570-pruning-reproduction}.
\end{abstract}

\section{Introduction}

Modern deep learning architectures routinely contain tens to hundreds of millions of
parameters, the majority of which are not necessary for inference-time performance
\citep{lecun1989optimal, han2015learning}.
Identifying and removing redundant weights---a process known as \emph{pruning}---is one
of the most widely studied model compression techniques, with applications in deploying
neural networks on resource-constrained hardware such as edge devices and mobile
processors.

The seminal work of \citet{han2015learning} demonstrated that fully connected and
convolutional networks can be compressed by 9--13$\times$ with no loss in top-5 accuracy
on ImageNet by (1) pruning weights below a magnitude threshold and (2) retraining the
sparse network.
A critical but sometimes overlooked implementation detail is the distinction between
\emph{one-shot} pruning---applying the full target sparsity in a single step---and
\emph{iterative} pruning, where the sparsity budget is gradually increased with fine-tuning
between rounds.

This reproduction project investigates the following central claim from Han et al.:

\begin{quote}
  \emph{Iterative pruning with fine-tuning achieves significantly higher accuracy than
  one-shot pruning, with the advantage concentrated at high sparsity levels ($>$70\%).}
\end{quote}

We deliberately restrict the experiment to a simple three-layer MLP on MNIST
so that (i) the full training-pruning pipeline can be executed on a single CPU or GPU
in under 30 minutes, (ii) results are easily reproducible and transparent, and
(iii) the qualitative trends from the original large-scale experiments can be validated
at minimal computational cost.
Beyond the direct replication, we extend the sparsity sweep to 95\% to locate the
``accuracy cliff'' for each method, compare global and layer-wise one-shot pruning strategies,
and analyze differential compressibility across the three network layers.

\section{Related Work}

\paragraph{Early pruning work.}
Optimal Brain Damage \citep{lecun1989optimal} and Optimal Brain Surgeon \citep{hassibi1993second}
are the classical approaches to weight removal based on second-order Taylor approximations
of the loss.
These methods are computationally expensive, motivating the magnitude-based heuristic
later popularized by \citet{han2015learning}.

\paragraph{Magnitude pruning and the lottery ticket hypothesis.}
\citet{han2015learning} demonstrated that globally thresholding by weight magnitude
yields competitive results after fine-tuning.
\citet{frankle2019lottery} subsequently showed that sparse sub-networks (``winning tickets'')
exist within the original network and can be re-trained from scratch to match dense
performance, reinforcing the idea that iterative pruning identifies structure that
the one-shot variant destroys.

\paragraph{Structured vs.\ unstructured pruning.}
Our work focuses on \emph{unstructured} (weight-level) pruning, where individual weights
are zeroed out.
\citet{wen2016learning} and subsequent work on structured pruning remove entire filters
or channels, yielding hardware-friendly sparse models but generally at higher accuracy cost
for a given compression ratio.

\paragraph{Fine-tuning and mask stability.}
A subtle but important implementation concern is that naïve fine-tuning after pruning
allows gradient updates to reactivate zeroed weights via optimizer momentum \citep{guo2016dynamic}.
Correct implementation requires masking gradients of pruned parameters to zero before
each optimizer step, which we implement explicitly.

\section{Methodology}

\subsection{Model and Dataset}

We use the MNIST handwritten-digit dataset \citep{lecun1998gradient}, comprising
60{,}000 training and 10{,}000 test images of $28{\times}28$ grayscale digits.
The model is a three-layer MLP with architecture $784 \to 300 \to 100 \to 10$
(ReLU activations, no dropout), matching the architecture studied in \citet{han2015learning}.
Training uses Adam \citep{kingma2015adam} with learning rate $10^{-3}$ and cross-entropy loss
for 10 epochs, achieving a dense baseline accuracy of 98.36\%.
All experiments use a fixed random seed for reproducibility.

\subsection{One-Shot Global Magnitude Pruning}

For a target sparsity $s \in [0, 1)$, we collect the magnitudes of all parameters in
the network into a single vector, compute the $s$-th quantile as the threshold $\tau$,
and set to zero all weights with $|w| < \tau$:
\begin{equation}
  \tau = \text{Quantile}_{s}\!\left(\{|w_i|\}_{i=1}^{N}\right), \quad
  w_i \leftarrow w_i \cdot \mathbf{1}[|w_i| \geq \tau].
\end{equation}
No fine-tuning is applied in the one-shot variant; accuracy is measured immediately
after pruning.

\subsection{Iterative Pruning with Fine-Tuning}

For a target sparsity $s$, we divide it into $K = \max(3, \lfloor s / 0.10 \rfloor)$
rounds of equal increments $\delta = s/K$.
After each pruning round we fine-tune for 3 epochs with the pruned-weight mask held fixed:
\begin{equation}
  \nabla_{w_i} \leftarrow \nabla_{w_i} \cdot \mathbf{1}[w_i \neq 0].
\end{equation}
This gradient masking prevents the Adam optimizer's first- and second-moment accumulators
from reactivating weights that have been pruned, ensuring that measured sparsity accurately
reflects model compression.

\subsection{Layer-Wise Pruning}

As a complementary strategy, we apply the same magnitude threshold independently within
each layer rather than globally.
This guarantees that every layer reaches exactly the target sparsity, whereas global pruning
may allocate different effective sparsity levels to different layers depending on their
weight magnitude distributions.

\section{Experiments}

\subsection{Main Result: Accuracy vs.\ Sparsity}

\begin{figure}[t]
  \centering
  \includegraphics[width=0.82\linewidth]{accuracy_vs_sparsity.pdf}
  \caption{Test accuracy as a function of sparsity for one-shot pruning (red) and
    iterative pruning with fine-tuning (green). The dashed horizontal line marks the
    dense baseline (98.36\%). The two methods diverge sharply beyond 70\% sparsity.}
  \label{fig:acc_sparsity}
\end{figure}

Figure~\ref{fig:acc_sparsity} shows the accuracy vs.\ sparsity curves for both methods
across sparsities from 0\% to 95\%, and Table~\ref{tab:results} records the numeric values.

\textbf{Moderate sparsity (0--70\%).}
Both methods maintain near-baseline accuracy with at most a 1.1 pp drop.
Iterative pruning actually achieves \emph{slightly above} dense accuracy at 70\% (98.31\%
vs.\ 98.36\%), a common artifact of the mild regularizing effect of pruning at moderate
sparsity levels.
This confirms the core claim from Han et al.\ that a large fraction of network weights
can be removed without meaningful performance degradation.

\textbf{High sparsity (80--90\%).}
The methods diverge dramatically.
One-shot pruning degrades from 95.21\% at 80\% to 80.43\% at 90\%---a sharp 14.78 pp
drop in a single step---while iterative pruning remains nearly flat at 98.19\% and 97.85\%
respectively.
The 17.42 pp gap at 90\% confirms that the fine-tuning recovery mechanism is not merely
helpful but essential at this compression regime.

\textbf{Extreme sparsity (95\%).}
The divergence widens catastrophically.
One-shot pruning collapses to 56.67\%---only marginally above a random classifier for
a 10-class problem---while iterative pruning retains 96.51\%, a 39.84 pp gap.
This result, which substantially extends the findings from our preliminary experiments,
demonstrates that one-shot pruning at extreme sparsity effectively destroys the network's
learned representations, whereas the iterative approach gradually restructures the
remaining weights to compensate for each pruning step.

\begin{table}[t]
  \centering
  \caption{Test accuracy (\%) at each sparsity level. Baseline accuracy is 98.36\%.
    The accuracy drop relative to the dense model is shown in parentheses.}
  \label{tab:results}
  \begin{tabular}{@{}lccc@{}}
    \toprule
    \textbf{Sparsity} & \textbf{One-Shot} & \textbf{Iterative + FT} & \textbf{Gap} \\
    \midrule
    0\%  (dense) & 98.36          & 98.36          & —       \\
    30\%         & 98.23 ($-$0.13) & 98.03 ($-$0.33) & 0.20 pp \\
    50\%         & 98.13 ($-$0.23) & 98.17 ($-$0.19) & 0.04 pp \\
    70\%         & 97.29 ($-$1.07) & 98.31 ($-$0.05) & 1.02 pp \\
    80\%         & 95.21 ($-$3.15) & 98.19 ($-$0.17) & 2.98 pp \\
    90\%         & 80.43 ($-$17.93) & 97.85 ($-$0.51) & 17.42 pp \\
    95\%         & 56.67 ($-$41.69) & 96.51 ($-$1.85) & 39.84 pp \\
    \bottomrule
  \end{tabular}
\end{table}

\subsection{Layer Compressibility Analysis}

\begin{figure}[H]
  \centering
  \includegraphics[width=0.78\linewidth]{layer_sparsity_heatmap.pdf}
  \caption{Per-layer sparsity achieved by iterative pruning targeting 90\% global
    sparsity, analogous to Figure~3 in \citet{han2015learning}.}
  \label{fig:heatmap}
\end{figure}

Figure~\ref{fig:heatmap} shows the actual proportion of zeroed weights per layer
after iterative pruning to a 90\% global target.
The first fully connected layer (fc1, $784 \to 300$, 235{,}200 parameters) accumulates
the most zeros under the global threshold, while the classification layer (fc3,
$100 \to 10$, 1{,}000 parameters) retains the most weights.
This differential compressibility arises naturally from global magnitude pruning:
the large first layer contains the most parameters and a correspondingly wide range of
weight magnitudes, so a global threshold removes proportionally more of its weights.
The small output layer, by contrast, cannot afford high sparsity---it directly determines
class scores from only 1{,}000 weights---and its weights are concentrated at higher
magnitudes as a result of training pressure on the final classification decision.

\subsection{Global vs.\ Layer-Wise Pruning}

At a 90\% sparsity target, the three one-shot and iterative strategies yield strikingly
different outcomes (Table~\ref{tab:layerwise}).

\begin{table}[h]
  \centering
  \caption{Comparison of pruning strategies at 90\% global sparsity.}
  \label{tab:layerwise}
  \begin{tabular}{@{}lcc@{}}
    \toprule
    \textbf{Strategy} & \textbf{Test Accuracy} & \textbf{Acc.\ Drop} \\
    \midrule
    Dense baseline           & 98.36\% & —        \\
    One-shot global          & 80.43\% & $-$17.93 pp \\
    One-shot layer-wise      & 42.88\% & $-$55.48 pp \\
    Iterative global + FT    & 97.85\% & $-$0.51 pp  \\
    \bottomrule
  \end{tabular}
\end{table}

Counter-intuitively, naive layer-wise pruning performs substantially \emph{worse} than
global pruning (42.88\% vs.\ 80.43\%).
The explanation lies in the architectural imbalance of the network: layer-wise pruning
at 90\% zeroes out 90\% of every layer uniformly, including the small output layer
(fc3, 1{,}000 weights), which is reduced to only 100 active weights spanning the full
$100 \to 10$ classification mapping.
Global pruning, by contrast, naturally concentrates zeros in the large first layer
(which has many redundant weights) and applies a lighter effective sparsity to the
critical output layer.
This result underscores a key practical insight: uniform per-layer sparsity budgets
are not appropriate when layers differ substantially in size, and sensitivity-aware
allocation is necessary for layer-wise strategies to be competitive.

\section{Conclusion}

We successfully reproduced and extended the central claim of \citet{han2015learning}
at small scale: magnitude-based pruning can remove the vast majority of network weights
with minimal accuracy loss, but only when fine-tuning is applied iteratively.
Our experiments on MNIST with a three-layer MLP confirm and sharpen the original findings:

\begin{enumerate}
  \item Both pruning methods retain near-baseline accuracy up to ${\sim}$70\% sparsity
    (within 1.1 pp of the dense model).
  \item One-shot pruning exhibits a sharp accuracy cliff between 80\% and 90\% sparsity,
    dropping 14.78 pp in a single step.
  \item At 95\% sparsity, iterative pruning (96.51\%) outperforms one-shot pruning
    (56.67\%) by 39.84 pp, demonstrating that fine-tuning is not merely beneficial
    but essential at extreme compression.
  \item Naive layer-wise pruning without fine-tuning (42.88\% at 90\%) is catastrophically
    worse than global pruning, revealing that uniform per-layer sparsity budgets
    over-compress critical small layers such as the output classifier.
\end{enumerate}

\textbf{Track-specific contribution (TinyReproductions).}
This work distills the core experimental story of Han et al.\ into a single reproducible
Python script runnable on a consumer laptop in under 15 minutes.
The extended sparsity sweep to 95\% and the layer-wise failure mode analysis provide
concrete insights not present in the original paper's small-scale experiments, and
the gradient-masking fine-tuning implementation correctly prevents optimizer momentum
from reactivating pruned weights---a common pitfall in reproduction efforts.

\bibliographystyle{iclr2026_conference}
\bibliography{references}

\section*{LLM Acknowledgement}

Claude Sonnet 4.6 (Anthropic) was used for assistance with LaTeX formatting and
Python code structure during the preparation of this report.
All experimental design, implementation logic, result analysis, and substantive
written content were produced by the author.

\end{document}
